\scanpage{127}
% archive sheet number: 126
% source lines 1--4
\deleted{où les $\handwritten{\sigma}_i$ (resp. $\handwritten{\sigma'}_i$) sont déduits de\deleted{s} $\handwritten{\sigma}$ (resp. de $\handwritten{\sigma'}$) par les trois homomorphismes $\handwritten{\delta}_i:U\overset{\handwritten{\scriptscriptstyle\wedge}}{\otimes}_kU\handwritten{\longrightarrow}U\overset{\handwritten{\scriptscriptstyle\wedge}}{\otimes}_kU\overset{\handwritten{\scriptscriptstyle\wedge}}{\otimes}_kU$ \deleted{utilisant} $\handwritten{\Delta}\otimes\mathrm{id}_U$ et symétriques.} (NB je n'ai pas écrit le calcul et ne garantis pas la forme exacte de la relation de l'hexagone).

% source lines 5--13
Soit $k'$ une algèbre (comm. ass. unitaire !) sur $k$, d'où $\unsure{\mathbb{U}|U}'=(U'_{\handwritten{\alpha}})$ et $\unsure{\mathbb{A}|A}'=\underset{\handwritten{\alpha}}{\underset{\handwritten{\longrightarrow}}{\lim}}A'_{\handwritten{\alpha}}$ sur $k'$, avec leurs structures de\deleted{\unsure{algèbres}} bigèbres \handwritten{(}resp. bi-progèbre pour $U'$\handwritten{)}. On se propose de déterminer les foncteurs $G:C\handwritten{\to}\deleted{\mathrm{Modtf}(}\mathrm{Projf}(k')$\inserted{(modules proj de t f sur $k'$)} \uline{exacts et} $k$-\uline{linéaires}, munis d'un isomorphisme $c:G(X\otimes Y)\handwritten{\simeq}G(X)\otimes G(Y)$. \deleted{\unsure{son}} Tout d'abord la donnée de $G^{\unsure{k}}$ exact à droite\inserted{$C\handwritten{\to}\mathrm{Mod}(k')$} revient à la donnée d'un $L=(L_{\handwritten{\alpha}})$, module\inserted{à droite} sur $U$\handwritten{,} muni d'une structure de $k'$-module \deleted{sur $k$} commutant aux opérations de $U$ (NB $L_{\handwritten{\alpha}}=G(U_{\handwritten{\alpha}\unsure{S|s}})$ ), par
\[
G(\mathbb{M})=L\otimes_U\mathbb{M}.
\]

% source lines 14--23
La $k^{\deleted{\unsure{'}}}$-linéarite signifie que les deux structures sur $k$ de $L$ coïncident, i.e. que $L$ provient d'un\deleted{\unsure{e}} module (pro) sur $U'=(U'_{\handwritten{\alpha}})$ (chaque $L_{\handwritten{\alpha}}$ est un module sur $U'_{\handwritten{\alpha}}$).\inserted{On peut alors écrire $G(X)=L\otimes_{U'}X'$, où $X'=X\otimes_kk'$.} Le fait que l'on s'envoie dans $\mathrm{Projf}(k')$ \deleted{s'exprime} et que le foncteur envisagé $G$ est exact s'exprime alors par le fait que les $L_{\handwritten{\alpha}}$ sont \deleted{libres}\inserted{proj.} de type fini \uline{sur} $k'$, et que les $L_{\handwritten{\unsure{\mu|\alpha}}}$ sont \unsure{plets|plats} sur \uline{les} $U_{\handwritten{\alpha}}$ (nécessité triviale, suffisance un petit argument par réduction au cas $X$ simple, puis $U$ un corps égal à $X$ \ldots). Donc on trouve exactement les \uline{modules} $L=(L_{\handwritten{\alpha}})$ \uline{sur} $U'$ \uline{tels que les} $\handwritten{L_\alpha}$ \uline{soient plats sur} $U_{\handwritten{\alpha}}$ \uline{et proj de tf sur} $k'$.
\margin{séparer les conditions, commencer par exactitude = platitude\handwritten{sur les $U_\alpha$}}
D'autre part, indépendamment de ces \deleted{C}\inserted{\handwritten{\unsure{c}}}onditions de platitude,\inserted{\handwritten{si dans $C$ il y a $1$ et $\vee$ \unsure{comp.tibles} avec\inserted{(\unsure{contravariété|contravariants})} $F$, i.e. $\Delta_U$ unitaire et $U$ avec antipodisme}} la donnée d'un isomorphisme $c$ revient à la donnée d'un isomorphisme de bi-$U'$-modules\inserted{\handwritten{à droite (*)}}
\margin{\deleted{\handwritten{\unsure{en le dualisant} \unsure{faisant intervenir} \illegible{2} d'un \unsure{antipodisme} $U'$}}}
\[
L\overset{\handwritten{\scriptscriptstyle\wedge}}{\otimes}_{\handwritten{k'}}L\handwritten{\xleftarrow{\sim}}L\overset{\handwritten{\scriptscriptstyle\wedge}}{\otimes}_{\handwritten{k'}}U'
\quad(\handwritten{\simeq}L\otimes_{(U,\handwritten{\Delta})}(U\overset{\handwritten{\scriptscriptstyle\wedge}}{\otimes}_kU)\deleted{\handwritten{\unsure{\simeq}}}\handwritten{\simeq}L\otimes_{(U',\handwritten{\Delta'})}(U'\overset{\handwritten{\scriptscriptstyle\wedge}}{\otimes}_kU')).
\]
\inserted{\handwritten{\unsure{[|(}\unsure{grâce|voir} à l'antipodisme:\unsure{)|]} $x\otimes u\longrightarrow x\otimes(u\otimes1)$]}}

% source lines 24--26 and lower-margin notes
\deleted{En termes} Utilisant les conditions de platitude, on peut définir $B=\mathrm{Homcont}_{k'}(L,k')$, qui est limite inductive des $B_{\handwritten{\alpha}}$ proj de tf sur $k'$ (à fortiori\inserted{\handwritten{il est}} plat sur $k'$), et une coargèbre sur $A'$ (limite de coalgèbres
\margin{\handwritten{Commençons \unsure{par} tout exprimer du côté de $U,L$, et ne dualisons qu'à la fin}}
\handwritten{(*) La loi de bim. dans $L\otimes_{\unsure{k'|U'}}U'$ est donnée par}
\[
\handwritten{\left\{\begin{aligned}
&\text{1ère loi : à droite des }\unsure{\text{mult.}}\text{ de }U'\\
&\text{2ème loi : donnée par celle de mod : droit}\\
&\text{sur }U\otimes U^o\quad[U^o\text{ opère à droite sur }U'\text{ !}]\\
&\text{via }(\mathrm{id}_U\otimes a.)(\Delta_U)
\end{aligned}\right.}
\]
